%------------------------------------------------------------------------------%
\begin{exercises}

%------------------------------------------------------------------------------%
\begin{exercise}
  Use o Teorema da Função Contínua para mostrar que
  \begin{mathlist}
  \item \lim_{n\to\infty} \sin  \left(\frac{\pi}{2}+\frac{1}{n}\right) = 1
  \item \lim_{n\to\infty} \ln   \left(1+\frac{1}{4^n}\right) = 0
  \item \lim_{n\to\infty} \tanh \left(\dfrac{\;n+\sfrac{1}{\,2^n\,}\;}{\;5-\sfrac{10}{\,5^n\,}\;}\right) = 1
  \end{mathlist}
\end{exercise}

% 23-56 - Stewart pt pg 633
%------------------------------------------------------------------------------%
\begin{exercise}
Verifique se a sequência converge ou não e calcule o limite das sequências convergentes.
\begin{mathlist}[2]
  \item a_n = e^{\sfrac{1}{n}}
  % 31 - 40
  \item a_n = \tan\left(\frac{2n\pi}{1+8n}\right)
  \item a_n = \sqrt{\frac{n+1}{9n+1}}
  \item a_n = \frac{n^2}{\sqrt{n^3+4n}}
  \item a_n = \exp\left(\frac{2n}{n+2}\right)
  \item a_n = \cos\left(\frac{n}{2}\right)
  \item a_n = \cos\left(\frac{2}{n}\right)
  \item a_n = \frac{\ln(n)}{\ln(2n)}
  \item a_n = \frac{e^n+e^{-n}}{e^{2n}-1}
  \item a_n = \frac{\arctan(n)}{n}
  % 41 - 50
  \item a_n = n^2 e^{-n}
  \item a_n = \ln(n\!+\!1)\! -\! \ln(n)
\end{mathlist}
\end{exercise}

%------------------------------------------------------------------------------%
%\begin{exercise}
%\begin{answer}
%\end{answer}
%\end{exercise}

%------------------------------------------------------------------------------%
%\begin{exercise}
%\begin{mathlist}[2]
%  \item
%  \item
%  \item
%\end{mathlist}
%\begin{answer}
%  \begin{alphalist}[3]
%    \item
%    \item
%    \item
%  \end{alphalist}
%\end{answer}
%\end{exercise}

\end{exercises}
%------------------------------------------------------------------------------%
